%% ThirdEye — MSR 2027 Technical Paper, full submission (10 pp + 2 pp refs)
%%
%% Build: Overleaf, or locally with
%%   latexmk -pdf main.tex
%% Requires the ACM acmart class (Overleaf: "ACM Conference Proceedings
%% Primary Article Template"). No local TeX install is assumed.
%%
%% DOUBLE-ANONYMOUS: MSR technical papers are reviewed double-anonymous.
%% For submission, add the `anonymous` option to \documentclass below and
%% replace the repository URL in Section 10 with an anonymised archive.
%% Names are left visible here because this draft doubles as a capstone
%% deliverable.

\documentclass[sigconf,nonacm]{acmart}

\usepackage{booktabs}
\usepackage{multirow}
\usepackage{microtype}

\newcommand{\ci}[2]{\,[#1,\,#2]}
\newcommand{\tool}{ThirdEye}

\begin{document}

\title{A False-Alarm Rate Is Not a Property of the Tool}
\subtitle{Four Unreported Choices in LLM Smart-Contract Audit Evaluation}

\author{Umar Ahmed}
\affiliation{%
  \institution{PES University}
  \city{Bengaluru}
  \country{India}}
\email{PES2UG23CS662@pesu.pes.edu}

\author{Anvita Agarwal}
\affiliation{%
  \institution{PES University}
  \city{Bengaluru}
  \country{India}}
\email{PES2UG23CS085@pesu.pes.edu}

\author{Tarun Basavaraju}
\affiliation{%
  \institution{PES University}
  \city{Bengaluru}
  \country{India}}
\email{PES2UG23CS644@pesu.pes.edu}

\author{V Lakshitha}
\affiliation{%
  \institution{PES University}
  \city{Bengaluru}
  \country{India}}
\email{PES2UG23CS665@pesu.pes.edu}

\begin{abstract}
Reported false-alarm rates for LLM-based security tools are read as properties of
the tool. We show that they are substantially properties of the benchmark, and
that the choices responsible are systematically unreported.

Recent work establishes that scoring \emph{pipelines} move benchmark outcomes in
cybersecurity question-answering by as much as eighty percentage points. We ask
the same question of a \emph{detection} task, which differs in a way that turns
out to matter: a detection benchmark must \textbf{construct a negative class} and
\textbf{define a matching rule}, and a question-answering benchmark does neither.
Those two constructions dominate the numbers.

We build a balanced benchmark of 2,250 Solidity contracts (1,125 safe :
1,125 vulnerable) whose safe class is stratified by \emph{how the safe label was
obtained}, and evaluate a model-diverse LLM council on 1,154 of them. Holding the
tool, the seed, the model digests and the machine fixed, four evaluation choices
each move the headline metric further than any tool-versus-tool difference we
found in the literature. \textbf{The same 1,154 verdicts, scored two defensible
ways, yield recall of 25.5\% or 82.2\%, and false-alarm rates of 15.7\% or
37.5\%.} Controlling for our own detector coverage, 31.6 points of the recall
spread remain attributable to the scoring protocol alone --- a choice that
changes no verdict at all.

We also report what this costs us. Scored type-matched, our recall falls inside
the band of the static analysers we set out to beat. Our original hypothesis,
that a model-diverse council raises precision, is refuted on the one early corpus
with a negative class. Adversarial arbitration, the standard precision mechanism,
yields no held-out gain. We release the benchmark, per-contract checkpoints for
every run, and a decision log recording each correction made during the work.
\end{abstract}

\keywords{smart contracts, benchmark construction, large language models,
evaluation methodology, false positives, reproducibility}

\maketitle

%% ===================================================================
\section{Introduction}
%% ===================================================================

A paper reporting a new LLM-based vulnerability detector reports a recall and a
false-alarm rate. A reader compares those two numbers against the previous
paper's two numbers and forms a view about which tool is better. That comparison
is close to meaningless, and the reason is not that the tools are similar.

Consider two figures from the literature, both verified against their source
papers. GPTScan~\cite{gptscan} reports a 4.39\% false-alarm rate on a curated set
of non-vulnerable contracts. Heimdallr~\cite{heimdallr} measures \emph{the same
tool} at 97.50\% on real audit-contest projects. Twenty-two-fold, with no change
to the tool. The difference lies entirely in which contracts were designated
``not vulnerable.''

Now consider the most-cited low-false-positive figure in the field:
Slither's 10.9\%~\cite{slither}. Read the denominator. It is the share of
\emph{flagged} contracts that manual review found wrong --- one minus precision
--- established by sampling at least fifty flagged contracts, on the reentrancy
detector alone, with no true-negative set at all. It is not $FP/(FP+TN)$ and it
cannot be converted into one from the published data. It is nonetheless the
number LLM tools are informally held against.

Berriche et al.~\cite{berriche} show that this is not peculiar to smart
contracts: across eight cybersecurity benchmarks and ten models, a single
scoring-pipeline choice can move a model's score by more than eighty percentage
points. Their taxonomy of fifteen failure modes covers stop-sequence handling,
extractor divergence, token budgets and aggregation order.

Every benchmark in that study is question-answering. \textbf{A multiple-choice
question cannot produce a false alarm.} Their failure modes are therefore defects
in the plumbing around a task whose ground truth is given. A detection benchmark
is a different object: it must \emph{construct} the class of things that are not
bugs, and it must \emph{decide} what it means to have found one. Neither
construction has a correct answer, both are usually unreported, and --- we find
--- both move the headline further than the tool does.

\paragraph{This paper.}
We build a detection benchmark designed so that those two constructions can be
varied and measured rather than assumed, and run four controlled experiments on
one corpus with one tool.

\begin{description}
\item[RQ1] Does the \emph{provenance of the safe label} change the false-alarm
  rate?
\item[RQ2] Does the \emph{matching rule} change recall?
\item[RQ3] Does what the model was \emph{permitted to read} change measured
  performance?
\item[RQ4] Does the \emph{execution environment} change the verdicts themselves?
\end{description}

\paragraph{Contributions.}
\begin{enumerate}
\item A balanced 2,250-contract Solidity benchmark whose negative class is
  stratified by label provenance --- to our knowledge the first that permits a
  false-alarm rate to be decomposed by how much the ``safe'' label can be
  trusted.
\item Four controlled measurements on one corpus, each isolating a single
  evaluation choice, with paired statistics where the design is paired.
\item A protocol-sensitivity result: identical verdicts scored any-slice versus
  type-matched give 82.2\% and 25.5\% recall. We decompose the 56.7-point spread
  into 31.6 points of protocol effect and ${\approx}25$ points of our own
  coverage gap, and report both separately rather than claiming the sum.
\item Seven silent failure modes observed during the work, six proposed
  invariants, and a documented instance of our own instrument violating one of
  them.
\item A released artifact: corpus, loaders, per-contract checkpoints, replay
  tooling, and a chronological decision log.
\end{enumerate}

%% ===================================================================
\section{Background and Related Work}
%% ===================================================================

\subsection{What the standard corpora can and cannot measure}

SmartBugs-curated~\cite{smartbugs} and the Web3Bugs contest
corpus~\cite{web3bugs} are the sets on which recall is most often reported. Both
are \textbf{all-positive}: every contract carries a known bug. Precision computed
on such a set is mechanically $1.000$ whenever recall is non-zero, and a
false-alarm rate is not defined at all. Our own early runs make the point
concretely --- on SmartBugs-curated our council records precision $1.000$ and
recall $0.881$, which says nothing whatever about false alarms.

Durieux et al.~\cite{durieux} is the largest evaluation in the field: nine tools
over 47,518 chain contracts and 69 annotated ones, 428,337 analyses, 564 days of
compute. They report that \textbf{93\% of the wild contracts are flagged
vulnerable}, and then --- correctly --- go no further than calling this evidence
of ``a considerable number of false positives,'' because the wild corpus is
unlabelled. Their annotated set, where the split would be possible, has 69
contracts and no safe class. One paper, a flag rate without a false-alarm rate
and recall without a negative class, not through any defect of method but because
no corpus offered both. That is the gap this benchmark exists to close.

\subsection{The claim we are \emph{not} making}

Berriche et al.~\cite{berriche} own the general observation that benchmark scores
are pipeline-dependent, and state it more strongly than our data can: swings
above eighty points across eight benchmarks and ten models. Heimdallr
\cite{heimdallr} owns the observation that LLM auditors carry catastrophic
real-world false-alarm rates. We take both as premises rather than findings.

Our claim is narrower and, we argue, unclaimed: in a \textbf{detection} task, the
construction of the negative class and the definition of the matching rule are
themselves experimental variables of first-order magnitude. Neither exists in the
QA benchmarks Berriche et al.\ audit. Their ``denominator inflation'' is a
counting error over attempted versus total questions; ours is a question about
what the denominator of a false-positive rate should \emph{be}.

\begin{table*}[t]
\centering
\caption{Positioning. Each row is verified against the source paper's own text.}
\label{tab:positioning}
\small
\begin{tabular}{@{}p{0.17\textwidth}p{0.38\textwidth}p{0.38\textwidth}@{}}
\toprule
\textbf{Prior work} & \textbf{What it establishes} & \textbf{What it leaves open} \\
\midrule
GPTScan~\cite{gptscan} &
Logic-vulnerability detection with GPT plus program analysis; a false-alarm rate
on a curated negative set &
The negative set is one curated slice; the 22$\times$ gap
to~\cite{heimdallr} is unexplained \\
Slither~\cite{slither} &
Fast, precise static analysis; 10.9\% ``false positives'' &
A different denominator ($1-$precision), one detector class, no true-negative
set \\
Durieux et al.~\cite{durieux} &
Nine tools at scale; 93\% of wild contracts flagged; 42\% combined detection &
The negative class is unlabelled, so a false-alarm rate is not computable \\
Heimdallr~\cite{heimdallr} &
LLM auditors at 97.5\% FPR on contest code &
Negatives are ``not in the report,'' i.e.\ absence of evidence \\
COMPSAC'26~\cite{compsac} &
Paired safe variants; local 4-bit FPR 34.09\%; lexical bias &
States ``no context-window analysis'' as its own limitation \\
Berriche et al.~\cite{berriche} &
Scoring pipelines move QA scores by more than 80 points &
Every benchmark is QA --- \textbf{no negative class exists} \\
\bottomrule
\end{tabular}
\end{table*}

%% ===================================================================
\section{Benchmark Construction}
%% ===================================================================

The benchmark is the instrument, so its construction is a contribution rather
than a preliminary.

\subsection{Balance and tiers}

2,250 Solidity contracts, balanced 1,125 safe : 1,125 vulnerable. Both halves are
stratified into three tiers. The safe tiers are ordered by \emph{how much the
safe label can be trusted} (Table~\ref{tab:tiers}).

\begin{table}[t]
\centering
\caption{Corpus tiers, with counts scored in the main run.}
\label{tab:tiers}
\small
\begin{tabular}{@{}lrp{0.42\columnwidth}@{}}
\toprule
\textbf{Tier} & \textbf{$n$} & \textbf{How the label was established} \\
\midrule
\multicolumn{3}{@{}l}{\emph{Safe}} \\
Audited libraries & 172 & OpenZeppelin / Solady --- heavily reviewed, widely
deployed \\
Audit-reviewed, clean & 240 & In a published audit's scope, no finding filed \\
Real-world, no report & 191 & Deployed, no public report --- \emph{absence of
evidence} \\
\midrule
\multicolumn{3}{@{}l}{\emph{Vulnerable}} \\
Curated & 214 & SmartBugs-style annotated vulnerabilities \\
Injected & 170 & A known bug deliberately introduced \\
Audit-report findings & 167 & A real finding from a published audit \\
\bottomrule
\end{tabular}
\end{table}

The safe tiers are the novel part. Tier~1 is as close to ground truth as this
domain offers; tier~3 is an inference from silence. A tool's false-alarm rate is
conventionally reported against whichever of these a paper happened to use.

\subsection{The negative class is a claim, not a fact}
\label{sec:labelnoise}

Because tier~3 is absence of evidence, some ``false alarms'' against it may be
real unreported bugs. We bound this by hand. A random sample (seed 2026) of 20 of
the 79 blocked safe contracts, plus one from an earlier read, was reviewed
against each contract's externally callable surface: 14 of 21 (70\%) are clear
tool errors, 4 (20\%) are unclear without a full read, and 2 (10\%) are possible
real bugs in a contract the datasets call safe.

So roughly 70\% of blocked safe contracts are genuine tool errors, and at most
30\% of the false-alarm rate could be label noise. The two candidate real defects
--- both with unmodified external \texttt{mint}/\texttt{withdraw} surfaces ---
are flagged as such in the artifact. This bound is self-assessed by the tool's
own authors and we treat it as a threat rather than a result
(\S\ref{sec:threats}).

\subsection{Sampling discipline}

Samples are \textbf{nested and seeded}: extending a run from $n$ to $n' > n$
yields a strict superset, so a longer run cannot silently change the population.
Every contract is checkpointed individually, which makes every table in this
paper re-scorable without re-running inference --- a property we rely on in
\S\ref{sec:rq2}, where a new result is produced from old checkpoints at zero
inference cost.

%% ===================================================================
\section{The Instrument}
%% ===================================================================

The tool under measurement is not the contribution, and we describe it only as
far as the experiments require.

\tool{} routes a contract to up to eight single-class specialists (reentrancy,
access control, arithmetic, business logic, oracle/price manipulation,
flash-loan/MEV, DoS/gas, proxy upgradeability), each pinned to a distinct
open-weights base model served locally by Ollama. A finding whose quoted evidence
does not appear verbatim in the source is discarded before aggregation.

\paragraph{Verdict rule.}
Surviving finding confidences are pooled as
$\mathit{risk} = 1 - \prod_i (1 - \mathit{conf}_i)$, blocking at $\tau = 0.925$.
This matters for reading our tables, because an earlier rule blocked on
\emph{any} surviving finding --- the same rule at $\tau \to 0$. Both appear in the
artifact and they are never interchanged (Table~\ref{tab:rule}).

\begin{table}[t]
\centering
\caption{The two rules on the same 1,154 verdicts. The shipped rule is what this
paper reports throughout.}
\label{tab:rule}
\small
\begin{tabular}{@{}lccccc@{}}
\toprule
\textbf{Rule} & \textbf{tp/fp/tn/fn} & \textbf{FPR} & \textbf{Recall} &
\textbf{Prec.} & \textbf{F1} \\
\midrule
OR-gate & 498/356/247/53 & 59.0\% & 90.4\% & 0.583 & 0.709 \\
\textbf{Noisy-OR} & \textbf{445/180/423/106} & \textbf{29.9\%} &
\textbf{80.8\%} & \textbf{0.712} & \textbf{0.757} \\
\bottomrule
\end{tabular}
\end{table}

The choice between the two is not a tuning detail, because the OR-gate is what a
naive council implementation produces and it is what our own earlier runs used.
Table~\ref{tab:splits} gives the comparison over 10 random dev/test splits
($n = 233$, 116 dev / 117 test, $\tau$ fitted on dev at $0.575$).

\begin{table}[t]
\centering
\caption{Verdict rule over 10 random dev/test splits, mean $\pm$ s.d. The
threshold wins 9 of 10 splits. A per-class weighted variant, tested as a control,
wins only 3 of 10 and is not reported as the mechanism.}
\label{tab:splits}
\small
\begin{tabular}{@{}lrrrr@{}}
\toprule
\textbf{Rule} & \textbf{FPR} & \textbf{Prec.} & \textbf{Recall} & \textbf{F1} \\
\midrule
OR-gate & 64.8\%\,$\pm$5.4 & 0.563 & 0.938 & 0.703 \\
\textbf{Noisy-OR} & \textbf{27.4\%}\,$\pm$7.8 & \textbf{0.733} & 0.824 &
\textbf{0.770} \\
\bottomrule
\end{tabular}
\end{table}

Per-class weighting beats a flat threshold on only 3 of 10 splits
(F1 $0.770$ against a threshold-only $0.790$), so we attribute the gain to
\textbf{thresholding, not weighting}, and report the simpler mechanism. We note
this because the weighted variant was the intuitively better design and was
implemented first; the control is the only reason it is not in the paper as a
contribution.

\paragraph{Where the council helps, and where it does not.}
On a balanced access-control slice (8 positive, 8 negative, seed 0) the council is
a clear win over a single LLM: recall $0.125 \rightarrow 0.750$, F1
$0.222 \rightarrow 0.857$, precision $1.000$ throughout. On Etherscan-50, the
only other early corpus with a negative class, \textbf{it loses}: precision
$0.067$ / recall $0.167$, against a single LLM's $0.182$ / $0.667$ and an
LLM+Slither pipeline's $0.140$ / $1.000$. That corpus has scraped, noisy ground
truth and we weight it accordingly --- but the project's original hypothesis was
that a model-diverse council raises precision, and on the available evidence that
hypothesis is \textbf{refuted}. We report it because a paper about honest
measurement that suppressed its own negative result would be self-refuting.

Adversarial LLM-as-judge arbitration, the standard precision mechanism, wins
\textbf{4 of 10} held-out splits --- no gain, despite visibly reducing the FPR
column. A reader shown only that column would have recorded it as a success.

%% ===================================================================
\section{Results}
%% ===================================================================

All experiments hold the tool, seed, model digests and (except in RQ4) the
machine fixed. Intervals are Wilson~\cite{wilson}; paired designs use
McNemar~\cite{mcnemar}. We never report Wald intervals, whose width collapses to
zero at the boundaries where several of these rates sit.

\subsection{RQ1 --- the provenance of the safe label}
\label{sec:rq1}

Same tool, same contracts, same rule. Only the origin of the safe label varies.

\begin{table}[t]
\centering
\caption{RQ1. False-alarm rate by how the safe label was obtained, with Wilson
intervals.}
\label{tab:rq1}
\small
\begin{tabular}{@{}lrrr@{}}
\toprule
\textbf{Safe tier} & \textbf{$n$} & \textbf{fp} &
\textbf{False-alarm rate} \\
\midrule
Audited libraries & 172 & 27 & \textbf{15.7\%}\ci{11.0}{21.9} \\
Audit-reviewed, clean & 240 & 90 & \textbf{37.5\%}\ci{31.6}{43.8} \\
Real-world, no report & 191 & 63 & \textbf{33.0\%}\ci{26.7}{39.9} \\
\bottomrule
\end{tabular}
\end{table}

The effect is \textbf{two-level, not a monotonic gradient}. Audited libraries
separate cleanly from both weaker tiers; the two weaker tiers overlap and are not
distinguishable from each other. An earlier version of this work described a
three-step gradient, which the intervals do not support.

\paragraph{Is this just code complexity?}
The obvious confound is that audited libraries might simply be smaller and
simpler than contest code, making the gradient a statement about contract size
rather than label provenance. It is not, and the data runs the other way. Median
contract size is 6,649 characters for audited libraries, against 3,531 for
audit-reviewed and 3,783 for real-world code. The tier with the \textbf{lowest}
false-alarm rate has contracts with nearly \textbf{twice the median size} of the
tier with the highest. If surface area drove false alarms we would observe the
opposite ordering. The size confound therefore works \emph{against} the measured
effect, which means the effect of label provenance is if anything
\textbf{understated} here.

The 21.8-point spread between tiers 1 and 2 is larger than the difference between
most published tools. A paper evaluating only against audited libraries would
report 15.7\%; the same tool evaluated against contest code reports 37.5\%.

\paragraph{The positive class is stratified too.}
The same construction question applies to the vulnerable half, and the answer is
milder but present. Detection rates by how the vulnerability was established are
86.4\%\ci{81.2}{90.4} on curated annotations, 89.4\%\ci{83.9}{93.2} on
deliberately injected bugs, and 70.7\%\ci{63.4}{77.0} on findings taken from
published audit reports. Injected bugs --- the cheapest positives to manufacture,
and therefore common in generated benchmarks --- are the easiest to find, and real
audit findings are the hardest, by an 18.7-point margin whose intervals do not
overlap. A benchmark built predominantly from injected bugs overstates recall in
the same way that a benchmark built predominantly from audited libraries
understates false alarms. We report the positive-class effect as secondary because
it does not, unlike RQ1 and RQ2, change the metric by more than the field's
tool-versus-tool differences.

\subsection{RQ2 --- what counts as a detection}
\label{sec:rq2}

Recall requires deciding when a flag \emph{counts}. Two protocols appear in the
literature and neither is usually named:

\begin{itemize}
\item \textbf{any-slice-positive} --- the tool flagged the contract, and the
  contract carries the label. This is what our \S4 table and most head-to-head
  comparisons use.
\item \textbf{type-matched} --- the finding that fired must itself name the
  labelled category. This is Durieux et al.'s unit, and GPTScan's.
\end{itemize}

Re-scoring the same checkpoints through the same shipped rule, changing nothing
but the protocol (denominator: the 349 vulnerable contracts carrying a mapped
DASP-10~\cite{dasp} label; 202 of the 551 carry none and are excluded),
any-slice-positive gives $287/349 = \mathbf{82.2\%}$ and type-matched gives
$89/349 = \mathbf{25.5\%}$.

\textbf{A 56.7-point spread with no change to any verdict} --- but that figure
mixes two different things, and we separate them rather than claim the sum. The
council has eight roles, five of which map unambiguously onto DASP-10, and
\textbf{158 of the 349 contracts (45.3\%) carry only classes no role covers}.
Part of the spread is therefore \emph{our} coverage gap, not a property of the
protocol. Restricting to the 187 contracts that carry \textbf{only} classes for
which a specialist exists --- so the tool was genuinely looking for the labelled
thing --- isolates the protocol effect (Table~\ref{tab:rq2}).

\begin{table}[t]
\centering
\caption{RQ2. Decomposing the protocol spread. The coverage-controlled row is the
paper's claim.}
\label{tab:rq2}
\small
\begin{tabular}{@{}lrrrr@{}}
\toprule
\textbf{Set} & \textbf{$n$} & \textbf{Any-slice} & \textbf{Type-matched} &
\textbf{Spread} \\
\midrule
All labelled & 349 & 82.2\% & 25.5\% & 56.7 pts \\
\textbf{Coverage-controlled} & \textbf{187} & \textbf{77.0\%} &
\textbf{45.5\%} & \textbf{31.6 pts} \\
\bottomrule
\end{tabular}
\end{table}

So roughly 25 points are attributable to our role set and \textbf{31.6 points to
the scoring protocol alone}. We take the controlled figure as the paper's claim.
It is still larger than any tool-versus-tool difference we found in the
literature, and it is obtained with the tool held fixed and every verdict
unchanged. The uncontrolled 56.7 points is not discarded --- it is the number a
reader of our headline recall would actually be misled by --- but it is reported
as a \emph{joint} effect of protocol and coverage, never as protocol alone.

\begin{table}[t]
\centering
\caption{RQ2 by vulnerability class. Our two largest labelled classes have no
specialist and are nonetheless ``detected'' at 89.0\% and 92.5\%.}
\label{tab:rq2class}
\small
\begin{tabular}{@{}lrrrc@{}}
\toprule
\textbf{DASP-10 class} & \textbf{$n$} & \textbf{Any-slice} &
\textbf{Type-matched} & \textbf{Spec.?} \\
\midrule
Unchecked low-level calls & 109 & 89.0\% & --- & \textbf{none} \\
Reentrancy & 86 & 76.7\% & 29.1\% & yes \\
Arithmetic & 79 & 77.2\% & 62.0\% & yes \\
Time manipulation & 40 & 92.5\% & --- & \textbf{none} \\
Access control & 17 & 88.2\% & 70.6\% & yes \\
Bad randomness & 9 & 66.7\% & --- & \textbf{none} \\
Denial of service & 8 & 62.5\% & 50.0\% & yes \\
Front running & 4 & 100.0\% & 0.0\% & yes \\
\bottomrule
\end{tabular}
\end{table}

Reentrancy --- the class the field has studied longest --- is flagged on 76.7\% of
its contracts but \emph{named} by the reentrancy specialist on 29.1\%. A separate,
earlier two-model council independently scored recall \textbf{zero} on bad
randomness (0/8), denial of service (0/6) and arithmetic (0/15), corroborating the
gap on a different run.

We report $[25.5\%, 82.2\%]$ as a bracket rather than either endpoint. 25.5\% is a
\emph{lower} bound: only five of eight roles have an unambiguous DASP slot, and we
deliberately refuse to stretch \texttt{business\_logic},
\texttt{oracle\_price\_manipulation} and \texttt{proxy\_upgradeability} into
``other,'' which would manufacture type-matched hits.

\paragraph{This makes comparison with prior work fair for the first time.}
Durieux et al.'s units are type-matched: nine tools combined reach 42\% of 115
tagged vulnerabilities, the best single tool 27\%, and Slither 17\%. Our 25.5\%
sits inside that band, while returning a verdict on 100\% of contracts against
Slither's 31\% (\S\ref{sec:calibration}). Reported as 82.2\%, the same comparison
would have read as a two-fold win over the entire field combined --- an artifact
of protocol, not of capability.

\subsection{RQ3 --- what the model was permitted to read}
\label{sec:rq3}

The local runtime truncates any prompt exceeding \texttt{num\_ctx}
\textbf{without error}, and abstentions are excluded from scoring. A
misconfiguration therefore \emph{removes rows} rather than lowering a score. 540
of 2,250 contracts (24.0\%) overflow a 4,096-token window, and the overflow is
class-skewed --- 350 vulnerable against 190 safe, because real buggy contracts are
larger.

Paired ablation on identical contracts, $4{,}096 \rightarrow 16{,}384$,
vulnerable class (36 paired, 20 scored in both arms): abstention
$44.4\% \rightarrow 0.0\%$, recall $12.5\% \rightarrow 81.2\%$, accuracy
$25.0\% \rightarrow 75.0\%$. McNemar $b = 2$, $c = 12$, $p = 0.016$. The larger
window is also \textbf{4.01$\times$ faster} (sign test $p < 0.0001$) ---
truncation triggers context shifting, so the misconfiguration cost both accuracy
and time while appearing in neither.

\paragraph{Limitation, stated rather than managed.}
The safe-class arm of this experiment is \emph{not} reported. Our run of it was
corrupted by a backend outage that left 145/176 and 108/176 rows as total council
failures; only 31 pairs scored in both arms, and those 31 are whatever ran while
the backend was alive. Consequently \textbf{RQ3 is established on the vulnerable
class only}, and our 29.9\% headline false-alarm rate may be an under-estimate by
an amount this work does not yet bound. We flag this as the most important open
item rather than reporting a figure from a compromised run.

\subsection{RQ4 --- which machine produced the verdicts}

Identical contract ids, seeds and byte-identical model digests, two hosts. On the
main arm ($n = 230$) the two hosts agree on \textbf{0.735} of verdicts, with
accuracy $0.722$ on one host and $0.622$ on the other. Roughly one verdict in four
flips across machines. A \texttt{num\_parallel = 1} control ($n = 227$, agreement
$0.727$) excludes batch parallelism as the cause (paired McNemar $p = 0.75$).
Backend and hardware nondeterminism is documented
elsewhere~\cite{pape}; what we add is its \textbf{propagation to security
verdicts}, and therefore to a published rate.

\paragraph{Capacity is not the lever (negative result).}
Restoring a 2.7$\times$ larger model on the three semantic roles --- a
one-variable change on the same 230 contracts --- \emph{increases} false alarms
$49 \rightarrow 74$ (McNemar $p = 0.00008$) while reducing misses
$15 \rightarrow 7$ ($p = 0.043$), for a net accuracy drop of
$0.722 \rightarrow 0.651$. Stronger associative reasoning reads more benign
patterns as vulnerabilities.

\subsection{Calibration against the literature}
\label{sec:calibration}

\begin{table}[t]
\centering
\caption{Every published rate we could verify, with its denominator printed
beside it. Two of these are not comparable to the others by construction.}
\label{tab:calibration}
\small
\begin{tabular}{@{}lrl@{}}
\toprule
\textbf{Source} & \textbf{Rate} & \textbf{Denominator} \\
\midrule
GPTScan, self-reported & 4.39\% & $FP/(FP{+}TN)$ \\
Slither, self-reported & 10.9\% & $1-$precision \\
\textbf{\tool{}, audited-library tier} & \textbf{15.7\%} & $FP/(FP{+}TN)$ \\
\textbf{\tool{}, overall} & \textbf{29.9\%} & $FP/(FP{+}TN)$ \\
COMPSAC'26, local 4-bit & 34.09\% & $FP/(FP{+}TN)$ \\
\textbf{\tool{}, audit-reviewed tier} & \textbf{37.5\%} & $FP/(FP{+}TN)$ \\
GPTScan, measured by~\cite{heimdallr} & 97.50\% & $FP/(FP{+}TN)$ \\
\bottomrule
\end{tabular}
\end{table}

We do not claim 29.9\% is good. We claim it is \emph{interpretable}: it carries a
stated negative class, a per-tier decomposition, a hand-audited label-noise
bound, and intervals.

\paragraph{Against Slither on 150 identical contracts.}
Table~\ref{tab:slither} gives the comparison on one contract set with one ground
truth, under the rule NO-GO iff at least one High or Medium finding.

\begin{table}[t]
\centering
\caption{Slither against \tool{} on 150 identical contracts. Coverage, not
precision, is the difference.}
\label{tab:slither}
\small
\begin{tabular}{@{}lrr@{}}
\toprule
& \textbf{Slither} & \textbf{\tool{}} \\
\midrule
Answered & \textbf{46 / 150} & \textbf{150 / 150} \\
Abstained & 104 (69\%) & 0 \\
tp / fp / tn / fn & 32 / 2 / 8 / 4 & --- \\
Precision where it answers & \textbf{0.941} & 0.712 \\
Recall where it answers & 0.889 & 0.808 \\
\bottomrule
\end{tabular}
\end{table}

Where Slither answers it is \textbf{more precise than we are} --- a 5.9\%
false-discovery rate ($2/34$) that independently corroborates the 10.9\% its own
paper reports. Its 104 abstentions were diagnosed with \texttt{solc} directly and
\textbf{exactly one} was attributable to our toolchain; the rest are multi-file
contest codebases and pragma mismatches, i.e.\ properties of the corpus rather
than of the tool. Slither is a high-precision, low-coverage instrument; we are a
moderate-precision, full-coverage one. The two are not ordered, and any single
scalar that claims to order them is hiding the abstention column.

\paragraph{Where the council is unambiguously the right design.}
On a balanced access-control slice with real negative signal (8 positive, 8
negative, seed 0), a single LLM finds 1 of 8 and the council finds 6 of 8, both at
precision $1.000$: recall $0.125 \rightarrow 0.750$, F1
$0.222 \rightarrow 0.857$. Adding arbitration on top returns it to $0.125$. This
is the clearest positive result in the work and it is obtained on 16 contracts,
which is why we present it as an existence proof for the mechanism rather than as
a performance claim.

\paragraph{Against GPTScan on 34 gradable shared projects.}
33/34 against 31/34, with \textbf{overlapping} intervals: no detection difference
is demonstrated, and we report that rather than a win.

%% ===================================================================
\section{Silent Failure Modes}
%% ===================================================================

Seven defects found during this work, each of which produced \emph{plausible}
metrics rather than obvious garbage. Directional bias is what makes them
dangerous.

\begin{table*}[t]
\centering
\caption{Seven silent failure modes observed during this work. In every case the
metric produced was plausible, and in six of seven the bias was towards a
\emph{better}-looking result.}
\label{tab:defects}
\small
\begin{tabular}{@{}rp{0.42\textwidth}p{0.42\textwidth}@{}}
\toprule
\textbf{\#} & \textbf{Defect} & \textbf{Effect on the reported metric} \\
\midrule
1 & Pinned model absent; at exactly 50\% errored, the ``$>$50\%'' fail-closed rule
does not trip &
A half-dead council records a clean GO. One 1,152-contract run invalidated \\
2 & Partial-council results treated as terminal &
32 of 198 rows had errored specialists, and \textbf{100\% of them were GO} \\
3 & Provider quota drained mid-run, checkpointed as terminal &
Bakes a transient outage permanently into recall \\
4 & Arbitration silently defaulted to a hosted backend &
The reported ``local tier'' was not local \\
5 & Arbiter config keyed on names the caller never passes &
Precision collapsed to $0.094$ and looked like a finding \\
6 & Subsampling by \texttt{items[:N]} &
Filename order clusters by project: the ``sample'' was 263 consecutive files \\
7 & Health-probe timeout shorter than a cold model load &
Aborts a healthy backend, so the run under-reports coverage \\
\bottomrule
\end{tabular}
\end{table*}

Table~\ref{tab:defects} lists them with the direction of the bias each induced.
Six of the seven make the tool look better than it is, which is the property that
makes them hard to notice: a result that improves is rarely audited as closely as
one that degrades.

\paragraph{Proposed invariants.}
(i)~A verdict is terminal only if the whole council ran. (ii)~Transient
$\neq$ deterministic: provider failures must never be checkpointed as results.
(iii)~Every stage records the backend and model it \emph{actually used}.
(iv)~Configuration lookups must fail loudly, never fall back. (v)~Samples must be
nested and seeded. (vi)~Probe timeouts must exceed measured worst-case cold
start.

\paragraph{We then violated invariant (ii) ourselves.}
The context-ablation runner --- the instrument measuring RQ3 --- wrote rows in
which every specialist had errored into its checkpoint directory, and its resume
path returned them verbatim, so a retry silently reused an outage instead of
repairing it. This is the failure reported in \S\ref{sec:rq3}. We record it
because it is the strongest available evidence that these invariants are not
obvious in retrospect: they were written by the authors, published in a prior
draft, and violated in the same codebase.

%% ===================================================================
\section{Threats to Validity}
\label{sec:threats}
%% ===================================================================

\paragraph{Construct.}
Tier-3 ``safe'' is absence of evidence. \S\ref{sec:labelnoise} bounds this at
${\approx}30\%$ by hand review, but the bound rests on 21 contracts assessed from
externally callable surfaces rather than full reads, and it was assessed by the
tool's own authors, unblinded. An independent second rater is the single cheapest
improvement available to this work.

\paragraph{Internal.}
RQ3 is established on the vulnerable class only (\S\ref{sec:rq3}). The
type-matched mapping in RQ2 covers five of eight roles; a more generous map would
raise 25.5\%, which is why it is reported as a lower bound and why the
coverage-controlled 31.6-point figure, not the 56.7-point one, is the claim. The
RQ1 effect is not explained by contract size, which runs opposite to it
(\S\ref{sec:rq1}), but we have not controlled for cyclomatic complexity or
external-call count, either of which could still carry part of it.

\paragraph{External.}
One tool, one language, one model family, two machines. The \emph{mechanisms} ---
constructed negatives, matching rules, silent truncation, execution
nondeterminism --- are not specific to any of these, but the magnitudes certainly
are.

\paragraph{Statistical.}
RQ3 rests on 20 paired contracts and RQ1's smallest tier on 172. We report Wilson
intervals throughout and never Wald.

%% ===================================================================
\section{What Is Not Claimed}
%% ===================================================================

\begin{itemize}
\item Not that \tool{} is a good auditor. A 29.9\% false-alarm rate is not
  deployable unattended, and scored type-matched our recall falls inside the band
  of the static analysers we set out to beat.
\item Not that the benchmark is representative of all Solidity.
\item Not that benchmark scores being configuration-dependent is a new
  observation --- Berriche et al.~\cite{berriche} established that, in a
  neighbouring domain, first.
\item Not that our rate is comparable to every figure in
  Table~\ref{tab:calibration}. Two are computed on a different denominator, which
  is why the denominator is printed beside each.
\item Not that we beat GPTScan. The head-to-head intervals overlap.
\end{itemize}

%% ===================================================================
\section{Conclusion}
%% ===================================================================

A false-alarm rate is not a property of a tool. In a detection task it is jointly
determined by how the negative class was constructed, what counts as a match,
what the model was permitted to read, and which machine ran it. We measured each
on one corpus with one tool, and each moved the headline further than the
differences the literature currently reports as results. The largest
well-controlled instance is 31.6 points, from a choice that changes no verdict at
all --- and even that is measured with the tool's own coverage held constant, so
it is a floor rather than a ceiling.

The practical consequence is a reporting standard, not a better tool. Concretely,
we propose that a paper reporting a detection rate on smart-contract code state
the following four things, none of which is expensive and none of which is
currently conventional.

\begin{enumerate}
\item \textbf{The provenance of the negative class.} Not its size --- its origin.
  ``Contracts with no filed audit finding'' and ``audited library code'' are
  different claims, and on our corpus they differ by 21.8 points with the tool
  unchanged. A paper that cannot say which it used has not reported a false-alarm
  rate; it has reported a number.
\item \textbf{The matching rule.} Whether a flag counts because the contract was
  flagged, or because the finding named the labelled class. On our data the two
  differ by 31.6 points with coverage controlled, and by 56.7 without. Both
  protocols are defensible; using one while being read as the other is not.
\item \textbf{The context window, and the abstention count beside it.} Silent
  truncation removes rows rather than lowering a score, so a truncated run reports
  a rate over a quietly different population. An abstention count makes this
  visible at a glance; a recall figure alone never does.
\item \textbf{The execution environment.} Host, backend, quantisation and batch
  settings, because one verdict in four flipped between two of our machines
  running byte-identical model digests.
\end{enumerate}

Absent those four, two reported rates are not comparable, and the field has been
comparing them anyway. We would rather this paper be cited for the checklist than
for the tool.

%% ===================================================================
\section{Data Availability}
%% ===================================================================

Corpus, loaders, per-contract checkpoints for every run, replay tooling, the
generated results index, and a chronological decision log recording each
correction made during this work are released at
\url{https://github.com/tarunb05/ThirdEye}. Every table above is re-scorable from
the committed checkpoints without re-running inference. Artifact scale: 18,078
lines of source across 61 commits, 42 backend regression tests pinning the
defects in \S6, and 3,707 lines of documentation.

\bibliographystyle{ACM-Reference-Format}
\bibliography{refs}

\end{document}
